\begin{textsample}{POS Dim 1 – vem\_ed\_26 – Score 39.00 – vem\_ed\_26/t141.txt}  \label{ex:f1_pos_059}
Artigo de \textbf{opinião} PCIs como política \textbf{linguística} na formação de tecnólogos No \textbf{contexto} atual, a globalização intensificou a \textbf{necessidade} de políticas públicas \textbf{que} promovam a inclusão e o acesso aos direitos fundamentais dos cidadãos – entre elas, as políticas \textbf{linguísticas}, \textbf{que} se \textbf{tornam} ferramentas estratégicas ao influenciar o \textbf{desenvolvimento} educacional, econômico e cultural das sociedades.
Segundo Calvet (2007, apud Dalla Corte, 2024), as políticas \textbf{linguísticas} abrangem decisões sobre seu uso social e refletem dinâmicas \textbf{institucionais}, econômicas e culturais.
Além disso, associam o ensino e a difusão de línguas ao planejamento público, incentivando a formação de cidadãos mais preparados para a comunicação internacional e para a \textbf{integração} cultural e econômica, como argumenta Martins (2016).
Este artigo discute a \textbf{importância} das políticas \textbf{linguísticas} e do ensino de \textbf{idiomas} estrangeiros, especialmente o espanhol, na formação \textbf{profissional} de tecnólogos em Comércio Exterior, \textbf{destacando} iniciativas como os Projetos Colaborativos Internacionais (PCIs / Cesu), \textbf{que} \textbf{fortalecem} a cooperação internacional e preparam os alunos para atuar de maneira eficaz em um \textbf{mercado} globalizado.
A \textbf{expansão} dos \textbf{mercados} internacionais reforçou a \textbf{importância} de políticas \textbf{linguísticas} \textbf{que} promovam a inclusão e o \textbf{desenvolvimento} social. \textbf{continuação} Elas emergem como um \textbf{elemento} estratégico ao integrar o \textbf{aprendizado} de \textbf{idiomas} ao \textbf{desenvolvimento} educacional, econômico e cultural.
O domínio de línguas estrangeiras, como o espanhol, exerce \textbf{papel} fundamental, especialmente nas relações com os países do Mercosul.
O Curso Superior de Tecnologia em Comércio Exterior oferecido nas Fatecs reflete essa \textbf{realidade} em seu Projeto Pedagógico de Curso (PPC), \textbf{que} possui quatro semestres de ensino do espanhol.
O objetivo \textbf{é} capacitar os alunos para a comunicação interpessoal e negociações, tornando-os mais preparados para atuarem em um cenário multicultural.
Nesse sentido, nota-se a \textbf{importância} de iniciativas de Internacionalização em Casa, como o PCI / Cesu realizado no primeiro semestre de 2024 entre alunos do 6º semestre de Comércio Exterior da Fatec Barueri e estudantes da DUOC UC (Chile), intitulado “Ponte Comercial: Praticando a Importação Bilateral Brasil-Chile”.
O PCI / Cesu reuniu alunos brasileiros e chilenos em grupos mistos, com objetivos como \textbf{promover} a \textbf{compreensão} \textbf{mútua}, desenvolver \textbf{habilidades} \textbf{práticas} e de liderança, além de aprimorar a comunicação em espanhol.
Durante três meses, o projeto \textbf{fortaleceu} laços de amizade e \textbf{consolidou} \textbf{conhecimentos} \textbf{práticos} entre os estudantes de ambos os países.
Essa experiência evidenciou o valor dos PCIs / Cesu, \textbf{que}, além de \textbf{promoverem} o \textbf{aprendizado} de uma língua estrangeira, aprofundam a \textbf{compreensão} intercultural necessária para negociações.
Considerando a relevância dos PCIs / Cesu, no segundo semestre de 2024, duas equipes de alunas, coautoras deste artigo, estão desenvolvendo \textbf{propostas} de Trabalho de Conclusão de Curso (TCC) com os temas “Língua Espanhola e Cultura nas Negociações Internacionais: Competências para o Egresso de Comércio Exterior” e “Políticas Linguísticas para a Cooperação entre Países: Um Estudo sobre o Mercosul no Contexto de Comércio Exterior”.
Essas pesquisas visam aprofundar a \textbf{compreensão} sobre como o domínio do espanhol pode \textbf{fortalecer} a comunicação e a diplomacia na região.
A proficiência em espanhol e as \textbf{habilidades} interculturais tornam-se, assim, diferenciais para os \textbf{futuros} tecnólogos de Comércio Exterior, \textbf{consolidando} o Brasil como um ator \textbf{relevante} no cenário econômico e cultural do Mercosul e de outros \textbf{mercados} internacionais. \textbf{Referências} DALLA CORTE, Anelise Copetti.
A Língua Espanhola no Contexto Universitário Sul-Brasileiro: uma análise das políticas \textbf{linguísticas} de internacionalização do espanhol. 2024. 200 f.
Tese (Programa de Pós-Graduação em Letras - Doutorado) - Universidade Estadual do Centro-Oeste, Guarapuava.
Disponível em: https://tede.unicentro.br/jspui/handle/jspui/2359.
Acesso em: 07 nov. 2024.
MARTINS, V.
Las \textbf{políticas} lingüísticas de enseñanza y difusión de español / lengua extranjera (ELE) en el Mercosur.
Onomázein Revista de Lingüística, Filología y Traducción, [ S. l. ], n. 33, p. 174-188, 5 ago. 2016.
DOI: 10.7764 / onomazein.33.10.
Disponível em: https://rda.uc.cl/index.php/onom/article/view/30073/23461.
Acesso em: 07 nov. 2024.

% matched lemmas: aprendizado, compreensão, conhecimento, consolidar, contexto, continuação, desenvolvimento, destacar, elemento, expansão, fortalecer, futuro, habilidade, idioma, importância, institucional, integração, linguístico, mercado, mútuo, necessidade, opinião, papel, político, profissional, promover, proposta, prático, que, realidade, referência, relevante, ser, tornar
\end{textsample}
